% GENERATED FILE. Edit canonical lessons, then run npm run generate:books.
% canonical-source-hash: fnv1a64:8fe5ada9319d1f6f
% canonical-lessons: AR-C305-mahkama, AR-C305-shahid-iyan, AR-C305-ittahama, AR-C305-muttaham, AR-C305-hakama

\chapter{In court}
\label{ch:ar-in-court}
\hlchaptermodality{305}

\begin{chapteropening}
\textbf{By the end of this chapter:} \emph{I can say a court, an eyewitness, to accuse, a defendant and to pass judgment.}

Complete the last lesson of chapter 305: I can say a court, an eyewitness, to accuse, a defendant and to pass judgment.
\end{chapteropening}

\section[\texorpdfstring{\ar{محكمة} (maḥkama)}{محكمة (maḥkama)}]{\ar{محكمة} (maḥkama) — a court of law}

\label{lesson:AR-C305-mahkama}



\begin{quote}
Before the new one: say the Arabic for at that time, then the Arabic for a reign.
\end{quote}

\subsection*{You'll want to know: \ar{محكمة}}
\textbf{\ar{محكمة}} — \emph{maḥkama} — \textquotedblleft{}a court of law\textquotedblright{}.

\begin{quote}
\ar{في} \ar{المدينة} \ar{محكمة} \ar{كبيرة}.
\end{quote}

\begin{quote}
\emph{fī al-madīna maḥkama kabīra.}
\end{quote}

\begin{quote}
There is a big court in the city.
\end{quote}

Plural: \ar{محاكم}, \emph{maḥākim}.

\begin{cousinweb}
\textbf{\ar{محكمة}} (\emph{maḥkama}) is built on the native Arabic root \emph{ḥ-k-m} (Lane), \textquotedblleft{}to restrain from wrong, to judge\textquotedblright{}. The \emph{mafʿala} pattern names a place: \textquotedblleft{}the place of judging\textquotedblright{}.
\end{cousinweb}

\begin{grammarlens}[title={What you've built}]
One more word for telling what happened.
\end{grammarlens}

\subsection*{Your turn}
\begin{itemize}
\raggedright
  \item \emph{Say it:} \emph{maḥkama}
  \item \emph{Say it:} \emph{maḥkama}, once more
  \item \emph{Say it:} \emph{fī al-madīna maḥkama kabīra.} — There is a big court in the city.
  \item \emph{Recall:} say the Arabic for at that time, then the Arabic for a reign, then say \emph{maḥkama} again
\end{itemize}

\begin{tcolorbox}[breakable,colback=teal!4,colframe=teal!35!black,title={Before you move on}]
What does \ar{محكمة} mean? (\textquotedblleft{}A court of law\textquotedblright{}.) Say it once more.
\end{tcolorbox}
\section[\texorpdfstring{\ar{شاهد} \ar{عيان} (shāhid ʿiyān)}{شاهد عيان (shāhid ʿiyān)}]{\ar{شاهد} \ar{عيان} (shāhid ʿiyān) — an eyewitness}

\label{lesson:AR-C305-shahid-iyan}



\begin{quote}
Before the new one: say the Arabic for a reign, then the Arabic for a court of law.
\end{quote}

\subsection*{You'll want to know: \ar{شاهد} \ar{عيان}}
\textbf{\ar{شاهد} \ar{عيان}} — \emph{shāhid ʿiyān} — \textquotedblleft{}an eyewitness\textquotedblright{}.

\begin{quote}
\ar{تكلم} \ar{شاهد} \ar{عيان} \ar{أمام} \ar{القاضي}.
\end{quote}

\begin{quote}
\emph{takallama shāhid ʿiyān amāma al-qāḍī.}
\end{quote}

\begin{quote}
An eyewitness spoke before the judge.
\end{quote}

Plural \ar{شهود} \ar{عيان}. The taught verb \ar{شاهد}, to watch, is the same root.

\begin{cousinweb}
\textbf{\ar{شاهد} \ar{عيان}} (\emph{shāhid ʿiyān}) is \textquotedblleft{}a witness by sight\textquotedblright{}: \emph{shāhid} is built on the native Arabic root \emph{sh-h-d}, \textquotedblleft{}to be present, to witness\textquotedblright{}, and \emph{ʿiyān} on \emph{ʿ-y-n}, the root of \emph{ʿayn}, an eye (Lane).
\end{cousinweb}

\begin{grammarlens}[title={What you've built}]
One more word for telling what happened.
\end{grammarlens}

\subsection*{Your turn}
\begin{itemize}
\raggedright
  \item \emph{Say it:} \emph{shāhid ʿiyān}
  \item \emph{Say it:} \emph{shāhid ʿiyān}, once more
  \item \emph{Say it:} \emph{takallama shāhid ʿiyān amāma al-qāḍī.} — An eyewitness spoke before the judge.
  \item \emph{Recall:} say the Arabic for a reign, then the Arabic for a court of law, then say \emph{shāhid ʿiyān} again
\end{itemize}

\begin{tcolorbox}[breakable,colback=teal!4,colframe=teal!35!black,title={Before you move on}]
What does \ar{شاهد} \ar{عيان} mean? (\textquotedblleft{}An eyewitness\textquotedblright{}.) Say it once more.
\end{tcolorbox}
\section[\texorpdfstring{\ar{اتهم} (ittahama)}{اتهم (ittahama)}]{\ar{اتهم} (ittahama) — to accuse}

\label{lesson:AR-C305-ittahama}



\begin{quote}
Before the new one: say the Arabic for a court of law, then the Arabic for an eyewitness.
\end{quote}

\subsection*{You'll want to know: \ar{اتهم}}
\textbf{\ar{اتهم}} — \emph{ittahama} — \textquotedblleft{}to accuse\textquotedblright{}.

\begin{quote}
\ar{اتهمت} \ar{الشرطة} \ar{الرجل} \ar{بالسرقة}.
\end{quote}

\begin{quote}
\emph{ittahamat ash-shurṭa ar-rajul bi-s-sariqa.}
\end{quote}

\begin{quote}
The police accused the man of theft.
\end{quote}

Present: \ar{يتهم}, \emph{yattahim}.

The person accused, then \ar{بـ} before the deed: \ar{اتهمه} \ar{بالكذب}.

\begin{cousinweb}
\textbf{\ar{اتهم}} (\emph{ittahama}) is built on the native Arabic root \emph{w-h-m} (Lane), \textquotedblleft{}to fancy, to suspect\textquotedblright{}. It is Form VIII, where \emph{w} merges into the \emph{t}: \textquotedblleft{}to cast suspicion on\textquotedblright{}; \emph{tuhma} is a charge.
\end{cousinweb}

\begin{grammarlens}[title={What you've built}]
One more word for telling what happened.
\end{grammarlens}

\subsection*{Your turn}
\begin{itemize}
\raggedright
  \item \emph{Say it:} \emph{ittahama}
  \item \emph{Say it:} \emph{ittahama}, once more
  \item \emph{Say it:} \emph{ittahamat ash-shurṭa ar-rajul bi-s-sariqa.} — The police accused the man of theft.
  \item \emph{Recall:} say the Arabic for a court of law, then the Arabic for an eyewitness, then say \emph{ittahama} again
\end{itemize}

\begin{tcolorbox}[breakable,colback=teal!4,colframe=teal!35!black,title={Before you move on}]
What does \ar{اتهم} mean? (\textquotedblleft{}To accuse\textquotedblright{}.) Say it once more.
\end{tcolorbox}
\section[\texorpdfstring{\ar{متهم} (muttaham)}{متهم (muttaham)}]{\ar{متهم} (muttaham) — a defendant, an accused}

\label{lesson:AR-C305-muttaham}



\begin{quote}
Before the new one: say the Arabic for an eyewitness, then the Arabic for to accuse.
\end{quote}

\subsection*{You'll want to know: \ar{متهم}}
\textbf{\ar{متهم}} — \emph{muttaham} — \textquotedblleft{}a defendant, an accused\textquotedblright{}.

\begin{quote}
\ar{قال} \ar{متهم} \ar{إنه} \ar{كان} \ar{في} \ar{بيته}.
\end{quote}

\begin{quote}
\emph{qāla muttaham innahu kāna fī baytihi.}
\end{quote}

\begin{quote}
One defendant said he had been at home.
\end{quote}

Plural \ar{متهمون}.

\begin{cousinweb}
\textbf{\ar{متهم}} (\emph{muttaham}) is built on the native Arabic root \emph{w-h-m} (Lane), \textquotedblleft{}to suspect\textquotedblright{}. \emph{Muttaham} is the passive participle of \emph{ittahama}: \textquotedblleft{}the one accused\textquotedblright{}.
\end{cousinweb}

\begin{grammarlens}[title={What you've built}]
One more word for telling what happened.
\end{grammarlens}

\subsection*{Your turn}
\begin{itemize}
\raggedright
  \item \emph{Say it:} \emph{muttaham}
  \item \emph{Say it:} \emph{muttaham}, once more
  \item \emph{Say it:} \emph{qāla muttaham innahu kāna fī baytihi.} — One defendant said he had been at home.
  \item \emph{Recall:} say the Arabic for an eyewitness, then the Arabic for to accuse, then say \emph{muttaham} again
\end{itemize}

\begin{tcolorbox}[breakable,colback=teal!4,colframe=teal!35!black,title={Before you move on}]
What does \ar{متهم} mean? (\textquotedblleft{}A defendant, an accused\textquotedblright{}.) Say it once more.
\end{tcolorbox}
\section[\texorpdfstring{\ar{حكم} (ḥakama)}{حكم (ḥakama)}]{\ar{حكم} (ḥakama) — to pass judgment; to rule}

\label{lesson:AR-C305-hakama}



\begin{quote}
Before the new one: say the Arabic for to accuse, then the Arabic for a defendant.
\end{quote}

\subsection*{You'll want to know: \ar{حكم}}
\textbf{\ar{حكم}} — \emph{ḥakama} — \textquotedblleft{}to pass judgment; to rule\textquotedblright{}.

\begin{quote}
\ar{حكم} \ar{القاضي} \ar{على} \ar{اللص} \ar{بالسجن}.
\end{quote}

\begin{quote}
\emph{ḥakama al-qāḍī ʿalā al-liṣṣ bi-s-sijn.}
\end{quote}

\begin{quote}
The judge sentenced the thief to prison.
\end{quote}

Present: \ar{يحكم}, \emph{yaḥkum}.

\ar{حكم} \ar{على} someone \ar{بـ}: to sentence; \ar{حكم} \ar{البلد}: to rule the country. The noun \ar{حكم} is a verdict.

\begin{cousinweb}
\textbf{\ar{حكم}} (\emph{ḥakama}) is built on the native Arabic root \emph{ḥ-k-m} (Lane), first \textquotedblleft{}to restrain from wrong\textquotedblright{}, then \textquotedblleft{}to judge\textquotedblright{} and \textquotedblleft{}to rule\textquotedblright{}; \emph{ḥikma}, wisdom, shares it.
\end{cousinweb}

\begin{grammarlens}[title={What you've built}]
One more word for telling what happened.
\end{grammarlens}

\subsection*{Your turn}
\begin{itemize}
\raggedright
  \item \emph{Say it:} \emph{ḥakama}
  \item \emph{Say it:} \emph{ḥakama}, once more
  \item \emph{Say it:} \emph{ḥakama al-qāḍī ʿalā al-liṣṣ bi-s-sijn.} — The judge sentenced the thief to prison.
  \item \emph{Recall:} say the Arabic for to accuse, then the Arabic for a defendant, then say \emph{ḥakama} again
\end{itemize}

\begin{tcolorbox}[breakable,colback=teal!4,colframe=teal!35!black,title={Before you move on}]
What does \ar{حكم} mean? (\textquotedblleft{}To pass judgment; to rule\textquotedblright{}.) Say it once more.
\end{tcolorbox}
